% WifiModule
\section{【応用】WifiModule — ステートマシンパターン}

\begin{patternbox}{WifiModule で学べるパターン}
\begin{itemize}
    \item enum classによる状態定義とswitch文による状態遷移
    \item ModuleTimerによるタイムアウト・再接続間隔の管理
    \item リクエストフラグによる外部からの制御（接続/切断指示）
    \item deinit()によるWiFiリソースの解放
\end{itemize}
\end{patternbox}

\subsection{状態定義}

\begin{lstlisting}[style=cppstyle, caption={WifiModule.h — 状態定義}]
enum class WifiState : uint8_t {
    DISCONNECTED,  // 未接続（接続リクエスト待ち）
    CONNECTING,    // 接続試行中
    CONNECTED,     // 接続済み
    RECONNECTING,  // 再接続試行中
    FAILED,        // 接続失敗（リトライ上限超過）
};
\end{lstlisting}

\begin{pointbox}[Data構造体にステートを公開]
\texttt{WifiData.state}としてenum値をSystemDataに公開することで、
ロジックフェーズが現在の接続状態に応じた判断を行える。
\texttt{isConnected}は簡易参照用のboolで、\texttt{state == CONNECTED}の省略形。
\end{pointbox}

\subsection{ステートマシンの実装}

\begin{lstlisting}[style=cppstyle, caption={WifiModule.cpp — updateInput()のステートマシン}]
void WifiModule::updateInput(SystemData& data) {
    // ステートマシン（WiFi状態の監視・自動再接続）
    switch (_state) {
        case WifiState::CONNECTING:
        case WifiState::RECONNECTING:
            if (WiFi.status() == WL_CONNECTED) {
                _state = WifiState::CONNECTED;
                _retryCount = 0;
            } else if (_stateTimer.getNowTime()
                       >= _config.connectTimeoutMs) {
                WiFi.disconnect();
                _retryCount++;
                if (_config.maxRetries > 0
                    && _retryCount >= _config.maxRetries) {
                    _state = WifiState::FAILED;
                } else {
                    _state = WifiState::RECONNECTING;
                    _stateTimer.setTime();
                }
            }
            break;

        case WifiState::CONNECTED:
            if (WiFi.status() != WL_CONNECTED) {
                _state = WifiState::RECONNECTING;
                _stateTimer.setTime();
            }
            break;
        // ...
    }

    // SystemDataに状態を反映
    data.wifi.state       = _state;
    data.wifi.isConnected = (_state == WifiState::CONNECTED);
    data.wifi.retryCount  = _retryCount;
}
\end{lstlisting}

\begin{lstlisting}[style=cppstyle, caption={WifiModule.cpp — updateOutput()のリクエスト処理}]
void WifiModule::updateOutput(SystemData& data) {
    // ロジックフェーズからのリクエスト処理
    if (data.wifi.requestConnect) {
        data.wifi.requestConnect = false;
        if (_state == WifiState::DISCONNECTED
            || _state == WifiState::FAILED) {
            _retryCount = 0;
            _startConnect();
        }
    }
    if (data.wifi.requestDisconnect) {
        data.wifi.requestDisconnect = false;
        _disconnect();
    }
}
\end{lstlisting}

\begin{pointbox}[ノンブロッキングなステートマシン]
WiFi接続は通常数秒かかるが、\texttt{delay()}で待たない。
\texttt{CONNECTING}状態では毎ループ\texttt{WiFi.status()}を確認し、
接続完了かタイムアウトかを判定する。
これにより、WiFi接続中も他のモジュール（LED、センサー等）は正常に動作し続ける。
\end{pointbox}

\begin{pointbox}[リクエストフラグパターン]
\texttt{requestConnect}はロジックフェーズからの「接続して」という指示。
BuzzerModuleの\texttt{requestPlay}と同じパターンで、
\texttt{updateOutput()}が処理後に\texttt{false}に戻す。
\end{pointbox}

\subsection{ロジックフェーズでの使用例}

\begin{lstlisting}[style=cppstyle, caption={ロジックフェーズでのWiFi制御}]
void applyPattern(SystemData& data) {
    // ボタンが押されたらWiFi接続開始
    if (data.button.justPressed
        && !data.wifi.isConnected) {
        data.wifi.requestConnect = true;
    }

    // WiFi接続状態に応じた表示
    switch (data.wifi.state) {
        case WifiState::CONNECTING:
            snprintf(data.display.line1, 48, "WiFi: Connecting...");
            break;
        case WifiState::CONNECTED:
            snprintf(data.display.line1, 48,
                     "WiFi: OK (RSSI=%d)", data.wifi.rssi);
            break;
        case WifiState::FAILED:
            snprintf(data.display.line1, 48,
                     "WiFi: Failed (%d retries)",
                     data.wifi.retryCount);
            break;
        default: break;
    }
}
\end{lstlisting}

% ─── API仕様 ───────────────────────────────
\subsection{API仕様}


\begin{apibox}{WifiModule --- WiFi接続管理モジュール}
\textbf{定義ファイル:} \texttt{lib/WifiModule/WifiModule.h / .cpp}\\
\textbf{分類:} 入出力モジュール（\texttt{updateInput()} + \texttt{updateOutput()}）\\
\textbf{対象デバイス:} ESP32 WiFi

WiFi接続のライフサイクル（接続・切断・再接続・タイムアウト）をステートマシンで自動管理するモジュール。
ロジックフェーズからリクエストフラグで接続・切断を制御する。
\end{apibox}

\subsubsection{機能概要}

\begin{itemize}
    \item 5状態のステートマシンによる接続ライフサイクル管理
    \item 自動再接続（リトライ回数制限付き）
    \item 接続タイムアウト検出
    \item RSSI（信号強度）とローカルIPの取得
    \item リクエストフラグによる接続・切断制御
    \item \texttt{deinit()}によるWiFi切断
\end{itemize}

% --- 列挙型 ---
\subsubsection{列挙型}

\begin{funcblock}{WifiState}
WiFi接続のステートマシン状態を定義する。

\begin{longtable}{l l p{6.0cm}}
\toprule
\textbf{値} & \textbf{名前} & \textbf{説明} \\
\midrule
\endhead
0 & \texttt{DISCONNECTED} & 未接続。\texttt{requestConnect}で接続開始 \\
1 & \texttt{CONNECTING}   & 接続試行中。タイムアウトまで待機 \\
2 & \texttt{CONNECTED}    & 接続済み。RSSI・IPを定期更新 \\
3 & \texttt{RECONNECTING} & 再接続試行中。切断検出後に自動遷移 \\
4 & \texttt{FAILED}       & 接続失敗。リトライ上限超過またはタイムアウト \\
\bottomrule
\end{longtable}
\end{funcblock}

% --- Config ---
\subsubsection{Config構造体}

\begin{funcblock}{WifiConfig}
WiFi接続のパラメータを定義する。

\begin{longtable}{l l l p{4.0cm}}
\toprule
\textbf{メンバ} & \textbf{型} & \textbf{制約・既定値} & \textbf{説明} \\
\midrule
\endhead
\texttt{ssid}                & \texttt{const char*}   & ---          & WiFi SSID \\
\texttt{password}            & \texttt{const char*}   & ---          & WiFiパスワード \\
\texttt{connectTimeoutMs}    & \texttt{unsigned long} & 例: 10000    & 接続タイムアウト [ms] \\
\texttt{reconnectIntervalMs} & \texttt{unsigned long} & 例: 5000     & 再接続までの待機時間 [ms] \\
\texttt{maxRetries}          & \texttt{uint8\_t}      & 例: 5（0=無制限）& 最大リトライ回数 \\
\bottomrule
\end{longtable}
\end{funcblock}

% --- Data ---
\subsubsection{Data構造体}

\begin{funcblock}{WifiData}
WiFi接続の状態と制御フラグを保持する。

\begin{longtable}{l l l p{3.8cm}}
\toprule
\textbf{メンバ} & \textbf{型} & \textbf{初期値} & \textbf{説明} \\
\midrule
\endfirsthead
\toprule
\textbf{メンバ} & \textbf{型} & \textbf{初期値} & \textbf{説明} \\
\midrule
\endhead
\midrule
\multicolumn{4}{r}{\small 次ページに続く} \\
\endfoot
\bottomrule
\endlastfoot
\multicolumn{4}{l}{\textbf{--- 状態（入力フェーズで更新） ---}} \\
\texttt{state}       & \texttt{WifiState} & \texttt{DISCONNECTED} & 現在のステートマシン状態 \\
\texttt{isConnected} & \texttt{bool}      & \texttt{false}        & 接続中フラグ（簡易参照用） \\
\texttt{retryCount}  & \texttt{uint8\_t}  & \texttt{0}            & 現在のリトライ回数 \\
\texttt{rssi}        & \texttt{int8\_t}   & \texttt{0}            & 信号強度 [dBm]（接続中のみ有効） \\
\texttt{localIp}     & \texttt{uint32\_t} & \texttt{0}            & ローカルIPアドレス（uint32\_t形式） \\
\midrule
\multicolumn{4}{l}{\textbf{--- 制御（ロジックフェーズで設定） ---}} \\
\texttt{requestConnect}    & \texttt{bool} & \texttt{false} & 接続リクエスト。\texttt{true}で接続開始 \\
\texttt{requestDisconnect} & \texttt{bool} & \texttt{false} & 切断リクエスト。\texttt{true}で切断 \\
\end{longtable}
\end{funcblock}

% --- メソッド ---
\subsubsection{メソッド}

\begin{funcblock}{WifiModule(const WifiConfig\& config)}
\begin{tabularx}{\textwidth}{l X}
\toprule
\textbf{項目} & \textbf{内容} \\
\midrule
引数   & \texttt{config} --- \texttt{WifiConfig}へのconst参照 \\
動作   & Configを\texttt{\_config}にコピー保持する \\
\bottomrule
\end{tabularx}
\end{funcblock}

\begin{funcblock}{bool init()}
\begin{tabularx}{\textwidth}{l X}
\toprule
\textbf{項目} & \textbf{内容} \\
\midrule
戻り値   & \texttt{bool} --- 常に\texttt{true} \\
動作     & WiFiモードをSTAに設定し、自動再接続を無効化する（ステートマシンで管理するため）。状態を\texttt{DISCONNECTED}に初期化 \\
\bottomrule
\end{tabularx}
\end{funcblock}

\begin{funcblock}{void updateInput(SystemData\& data)}
\begin{tabularx}{\textwidth}{l X}
\toprule
\textbf{項目} & \textbf{内容} \\
\midrule
書き込み先     & \texttt{data.wifi} \\
タイミング制御 & 毎ループ実行 \\
動作           & WiFi接続状態を確認し、ステートマシンの状態遷移を処理。\texttt{state}・\texttt{isConnected}・\texttt{rssi}・\texttt{localIp}・\texttt{retryCount}を更新 \\
\bottomrule
\end{tabularx}
\end{funcblock}

\begin{funcblock}{void updateOutput(SystemData\& data)}
\begin{tabularx}{\textwidth}{l X}
\toprule
\textbf{項目} & \textbf{内容} \\
\midrule
読み取り元     & \texttt{data.wifi} \\
タイミング制御 & 毎ループ実行 \\
動作           & \texttt{requestConnect}が\texttt{true}なら接続開始、\texttt{requestDisconnect}が\texttt{true}なら切断。各フラグを処理後に\texttt{false}にクリア \\
\bottomrule
\end{tabularx}
\end{funcblock}

\begin{funcblock}{void deinit()}
\begin{tabularx}{\textwidth}{l X}
\toprule
\textbf{項目} & \textbf{内容} \\
\midrule
動作   & WiFi接続を切断し、WiFiモードをOFFにする \\
\bottomrule
\end{tabularx}
\end{funcblock}

\begin{notebox}[ステートマシン遷移]
状態遷移の概要:

\begin{itemize}
    \item \texttt{DISCONNECTED} $\rightarrow$ \texttt{CONNECTING}（\texttt{requestConnect}）
    \item \texttt{CONNECTING} $\rightarrow$ \texttt{CONNECTED}（接続成功）
    \item \texttt{CONNECTING} $\rightarrow$ \texttt{FAILED}（タイムアウト）
    \item \texttt{CONNECTED} $\rightarrow$ \texttt{RECONNECTING}（切断検出）
    \item \texttt{CONNECTED} $\rightarrow$ \texttt{DISCONNECTED}（\texttt{requestDisconnect}）
    \item \texttt{RECONNECTING} $\rightarrow$ \texttt{CONNECTED}（接続成功）
    \item \texttt{RECONNECTING} $\rightarrow$ \texttt{FAILED}（リトライ超過）
    \item \texttt{FAILED} $\rightarrow$ \texttt{CONNECTING}（\texttt{requestConnect}、リトライカウントリセット）
\end{itemize}
\end{notebox}

% --- 使用例 ---
\subsubsection{使用例}

\begin{lstlisting}[style=cppstyle, caption={WifiModuleの使用例}]
// ProjectConfig.h
const WifiConfig WIFI_CONFIG = {
    .ssid                = "YourSSID",
    .password            = "YourPassword",
    .connectTimeoutMs    = 10000,  // 10秒タイムアウト
    .reconnectIntervalMs = 5000,   // 5秒後に再接続
    .maxRetries          = 5,      // 最大5回リトライ
};

// ロジックフェーズでの使用
void applyPattern(SystemData& data) {
    // ボタンでWiFi接続トグル
    if (data.button.justPressed) {
        if (data.wifi.isConnected) {
            data.wifi.requestDisconnect = true;
        } else {
            data.wifi.requestConnect = true;
        }
    }
    // 接続中はRSSIを表示
    if (data.wifi.isConnected) {
        snprintf(data.display.line3, 48,
                 "RSSI: %d dBm", data.wifi.rssi);
    }
}
\end{lstlisting}
